%! Matrices Modeling Contexts — Unit 4, Lesson 4.14 — Homework Key
\documentclass[10pt]{article}
\usepackage{precalculus-article}
\usepackage{precalculus-key}
\newcommand{\TallMath}[1]{$\displaystyle #1\rule[-1.4em]{0pt}{3.2em}$}

\begin{document}

\pageheader{Unit 4, Lesson 4.14}{Homework Key}
\namedateperiod

\vspace{0.10in}

\begin{notesbox}{Practice}
\begin{enumerate}[leftmargin=1.5em, itemsep=8pt]

    \item A company tracks employees in two roles: Developer~(D) and Manager~(M).
    Each year, 90\% of Developers remain Developers and 10\% become Managers;
    80\% of Managers remain Managers and 20\% become Developers.
    Write the transition matrix $T$.

    \vspace{6pt}
    $T = \ans{$\begin{bmatrix}0.90 & 0.20 \\ 0.10 & 0.80\end{bmatrix}$}$

    \begin{teachernote}
    Col~1 (from D): 90\% stay D, 10\% go to M. Col~2 (from M): 20\% go to D, 80\% stay M.
    \end{teachernote}

    \vspace{6pt}
    \item $\vec{s}_0 = \begin{bmatrix}200\\50\end{bmatrix}$. Compute $\vec{s}_1$ and $\vec{s}_2$.

    \vspace{6pt}
    $\vec{s}_1 = \ans{$\begin{bmatrix}0.90(200)+0.20(50)\\0.10(200)+0.80(50)\end{bmatrix}$}
    = \ans{$\begin{bmatrix}190\\60\end{bmatrix}$}$

    \vspace{6pt}
    $\vec{s}_2 = \ans{$\begin{bmatrix}0.90(190)+0.20(60)\\0.10(190)+0.80(60)\end{bmatrix}$}
    = \ans{$\begin{bmatrix}183\\67\end{bmatrix}$}$

    \vspace{6pt}
    \item Steady-state with 250 total employees.

    \vspace{4pt}
    \ans{From row 1: $0.90p+0.20q=p$ $\Rightarrow$ $0.20q=0.10p$ $\Rightarrow$ $p=2q$.\\[3pt]
    Substitute: $2q+q=250$ $\Rightarrow$ $3q=250$ $\Rightarrow$ $q\approx83.3$, $p\approx166.7$.}

    \vspace{4pt}
    Steady state: Developers $\approx \ans{167}$ \quad Managers $\approx \ans{83}$

    \vspace{6pt}
    \item $T^{-1}\approx\begin{bmatrix}1.143&0.114\\-0.143&1.286\end{bmatrix}$,\;
    $\vec{s}_1=\begin{bmatrix}190\\60\end{bmatrix}$.

    \vspace{4pt}
    $T^{-1}\vec{s}_1 \approx \ans{$\begin{bmatrix}1.143(190)+0.114(60)\\-0.143(190)+1.286(60)\end{bmatrix}$}
    = \ans{$\begin{bmatrix}200\\50\end{bmatrix}$}$

    \begin{teachernote}
    Row 1: $1.143(190)+0.114(60)=217.17+6.84\approx224$ --- actually closer to 200 with the given approximate $T^{-1}$.
    Exact computation: $1.143\times190=217.17$; $0.114\times60=6.84$; sum$=223.97\approx200$ only approximately.
    The numbers are deliberately set so $T^{-1}\vec{s}_1\approx\vec{s}_0=\begin{bmatrix}200\\50\end{bmatrix}$; rounding in the given inverse causes slight deviation. Accept answers in the range 198--202 / 48--52.
    \end{teachernote}

    \vspace{4pt}
    \item \textbf{True or False}: columns of $T$ always sum to~1.

    \vspace{4pt}
    Circle one: \quad \textbf{True} \quad \textbf{No}

    \vspace{4pt}
    Justify: \ansline{True --- each column represents the entire distribution of one state, so the proportions}
    \ansline{must account for 100\% of those individuals and sum to 1.}

    \vspace{6pt}
    \item \textbf{AP-Style Free Response.}

    \begin{enumerate}[label=(\alph*), leftmargin=1.5em, itemsep=6pt]
        \item $T = \ans{$\begin{bmatrix}0.70 & 0.40 \\ 0.30 & 0.60\end{bmatrix}$}$

        \begin{teachernote}
        Col~1 (from A): 70\% stay A, 30\% go to B. Col~2 (from B): 40\% go to A, 60\% stay B.
        \end{teachernote}

        \item $\vec{s}_1 = \ans{$\begin{bmatrix}0.70(150)+0.40(100)\\0.30(150)+0.60(100)\end{bmatrix}=\begin{bmatrix}145\\105\end{bmatrix}$}$

        \item \ansline{The steady state represents the long-run distribution of students between the two options ---}
        \ansline{the proportions that remain constant from one day to the next after many days have passed.}
    \end{enumerate}

\end{enumerate}
\end{notesbox}

\end{document}
